\documentclass[a4paper,11pt]{article}

\usepackage[greek,english]{babel}
\usepackage[utf8]{inputenc}
\usepackage[T1]{fontenc}
\newcommand{\gr}[1]{\foreignlanguage{greek}{#1}}

\usepackage{amsmath,amssymb,graphicx,booktabs,float,subcaption}
\usepackage{hyperref,listings,xcolor,multirow,array,enumitem}
\usepackage[left=2.5cm,right=2.5cm,top=2.5cm,bottom=2.5cm]{geometry}

\lstdefinestyle{mystyle}{
    backgroundcolor=\color{white},
    basicstyle=\ttfamily\footnotesize,
    breaklines=true, captionpos=b, keepspaces=true,
    numbers=left, numbersep=5pt, tabsize=2
}
\lstset{style=mystyle}

\title{Graph Neural Networks for Predicting Drug Safety Properties\\[0.3em]
\large Mid-Term Progress Report}
\author{
\gr{Κωνσταντίνα Πάσχου} \and
\gr{Γεώργιος Κεσογλίδης} \and
\gr{Θεοχάρης Αγγελίδης}\\[0.4em]
\small Course: Advanced Topics in Machine Learning\\
\small MSc Artificial Intelligence --- Aristotle University of Thessaloniki
}
\date{April 2026}

\begin{document}
\maketitle

\begin{abstract}
We investigated whether Graph Neural Networks provide an advantage over traditional fingerprint-based methods for drug toxicity prediction. We tested three representation strategies (Morgan fingerprints, RDKit descriptors, ChemBERTa) and four GNN architectures (GCN, GIN, GAT, GraphSAGE) across six toxicity datasets and twelve Tox21 endpoints (17 tasks total). On balanced datasets, classical ML with fingerprints and GNNs achieve comparable performance (0.79--0.92 AUROC). On severely imbalanced datasets (ClinTox 8\% positive, Tox21 3--15\%), all methods plateau around 0.79 AUROC for fingerprints and 0.54 AUROC for GNNs. Regression on LD50 shows fingerprints ($R^2 = 0.60$) outperform GNNs ($R^2 = 0.49$). The central finding: representation choice does not determine performance on imbalanced toxicity datasets. Class imbalance itself emerges as the limiting factor, not algorithm design. The next phase tests whether combining data-side interventions (rejection sampling) with GNNs can break this imbalance ceiling.
\end{abstract}

\tableofcontents
\newpage

% ─────────────────────────────────────────────────────────────────────────────
\section{Introduction and Scope}

Drug development is expensive and slow. A typical small-molecule program costs \$1--2 billion and takes 10--15 years, with over 90\% of candidates failing in clinical trials. Predicting toxicity computationally before synthesis could substantially reduce both timelines and costs.

Molecular property prediction has traditionally used hand-engineered descriptors such as LogP and TPSA, or structural fingerprints like Morgan/ECFP~\cite{chemrev}. Recently, Graph Neural Networks~\cite{kipf2017,xu2019,velickovic2018,hamilton2017} have shown promise on other molecular tasks. The central question we investigate is whether representation choice matters when the real constraint is data imbalance. Most real toxicity datasets are severely imbalanced: 8\% toxic in ClinTox, 3--15\% positive in Tox21 endpoints. Do better representations (learned embeddings, graphs) overcome class imbalance, or does imbalance itself become the ceiling regardless of representation?

We investigate this question across six toxicity datasets and twelve Tox21 endpoints through three parallel experimental investigations.

% ─────────────────────────────────────────────────────────────────────────────
\section{Datasets and Experimental Design}

\subsection{Domain Context: Drug Toxicity Prediction}

Drug toxicity prediction is fundamentally a class imbalance problem. A typical toxicity dataset contains thousands of non-toxic molecules and a small set of toxic ones. For example, ClinTox has 1,130 molecules with only 92 toxic (8\% positive), and Tox21 endpoints range from 3--15\% positive. This imbalance challenges all machine learning methods: a model that always predicts "non-toxic" achieves 92\% accuracy on ClinTox despite being useless.

\subsection{Molecular Representations}

All methods begin with SMILES strings, a text notation for molecular structure. The choice of representation fundamentally determines what information the model can access.

Fingerprints (Morgan/ECFP4) compress molecules into fixed-length binary vectors. Each bit represents whether a specific substructure pattern appears. With 1024 bits and radius 2, they capture local structural patterns. RDKit 2D descriptors encode hand-engineered scalar properties: molecular weight, LogP, TPSA, H-bond counts. These fixed-vector representations have dominated toxicity prediction for decades.

Graph representations preserve molecular structure as explicit nodes (atoms) and edges (bonds). SMILES are parsed into molecular graphs where nodes carry learned features (atomic number, formal charge, H-count, hybridization, aromaticity, ring membership, valence, radical electrons). Graphs are variable-size: the same model applies to molecules with 10 atoms or 100 atoms. Graph Neural Networks learn how molecular properties relate to toxicity directly from data.

Pre-trained language models (ChemBERTa) treat SMILES as text sequences. ChemBERTa uses a RoBERTa backbone pre-trained on 77 million SMILES strings, then fine-tuned for specific toxicity tasks. This approach combines the benefits of large-scale pre-training with task-specific adaptation.

\subsection{Evaluation Protocol}

All experiments use scaffold-based splits provided by Therapeutics Data Commons. Scaffold splits group molecules by their core chemical structure, ensuring test molecules are structurally different from training examples. This prevents inflated performance from similar molecules leaking across train/test splits.

Standard metrics are used throughout: ROC-AUC (threshold-independent, handles imbalance correctly), F1 (shows whether models predict the minority class), and R² (for regression tasks). Models are evaluated on fixed 0.5 threshold and calibrated thresholds (set to the positive-class ratio, e.g., 0.08 for ClinTox). This distinction matters: fixed 0.5 often causes F1 collapse on imbalanced data, while calibrated thresholds recover it.

\subsection{Datasets}

This work evaluates on six toxicity datasets from Therapeutics Data Commons (TDC)~\cite{huang2021} plus twelve Tox21 binary classification endpoints (17 tasks total):

\begin{itemize}
    \item \textbf{AMES} (Salmonella mutagenicity): 7,255 molecules, 50\% positive. Balanced classification.
    \item \textbf{hERG\_Karim} (cardiac arrhythmia risk): 13,445 molecules, 40\% positive. Moderate class balance.
    \item \textbf{ClinTox} (clinical trial outcomes): 1,484 molecules, 8\% positive. Severe imbalance.
    \item \textbf{DILI} (drug-induced liver injury): 475 molecules, 45\% positive. Balanced but small.
    \item \textbf{LD50\_Zhu} (acute systemic toxicity): 7,385 molecules, continuous regression target. No class labels.
    \item \textbf{Tox21} (nuclear receptor / stress-response assays): 7,831 molecules, 12 binary endpoints, 3--15\% positive each.
\end{itemize}

% ─────────────────────────────────────────────────────────────────────────────
\section{Results and Analysis}

\subsection{Baseline Performance: Classical ML with Fingerprints}

We establish the performance ceiling for hand-engineered features by testing Random Forest and XGBoost with Morgan fingerprints (ECFP4) on all six toxicity datasets. These results answer the first question: what is the benchmark any learned representation must beat?

\begin{table}[H]
\centering
\caption{Classical ML with Morgan fingerprints: Random Forest and XGBoost across all six toxicity datasets.}
\label{tab:baseline_classical}
\begin{tabular}{lrr}
\toprule
\textbf{Dataset} & \textbf{RF AUROC} & \textbf{XGB AUROC} \\
\midrule
AMES (50\% pos) & 0.906 & 0.881 \\
hERG\_Karim (40\% pos) & 0.920 & 0.897 \\
ClinTox (8\% pos) & 0.792 & 0.870 \\
DILI (45\% pos) & 0.868 & — \\
Tox21 (macro, 3--15\%) & 0.776 & 0.803 \\
\midrule
\multicolumn{3}{c}{\textbf{Regression}} \\
\midrule
LD50\_Zhu & \multicolumn{2}{c}{$R^2 = 0.58$ (RF), $R^2 = 0.60$ (XGB)} \\
\bottomrule
\end{tabular}
\end{table}

On balanced datasets (AMES, hERG), fingerprints achieve 0.91--0.92 AUROC. On severe imbalance (ClinTox 8\%), performance drops to 0.79 AUROC. On regression (LD50), classical ML plateaus at $R^2 = 0.60$. These numbers establish the fingerprint ceiling. Any other representation must beat these benchmarks on at least one metric to claim an advantage.

\subsection{Comparative Analysis of Molecular Representations}

To explore whether learned representations outperform hand-engineered features, three representation strategies were compared on AMES (balanced, 50\% positive) using scaffold-based splits. AMES provides a clean test case: abundant data, balanced classes, and high chemical diversity. Three representations were tested: Morgan fingerprints (ECFP4, 1024-bit, radius 2), RDKit 2D descriptors (MW, LogP, TPSA, H-bond donors/acceptors), and ChemBERTa (RoBERTa backbone pre-trained on 77M SMILES, fine-tuned with frozen backbone and trainable head).

\begin{table}[H]
\centering
\caption{Representation comparison on AMES (scaffold split).}
\label{tab:repr_ames}
\begin{tabular}{llrrr}
\toprule
\textbf{Representation} & \textbf{Model} & \textbf{AUROC} & \textbf{F1} & \textbf{Accuracy} \\
\midrule
Morgan FP (1024-bit) & RF & 0.900 & 0.839 & 0.825 \\
Morgan FP (1024-bit) & XGB & 0.899 & 0.842 & 0.826 \\
RDKit 2D descriptors & RF & 0.867 & 0.806 & 0.786 \\
ChemBERTa (77M pre-trained) & frozen backbone + head & 0.802 & 0.750 & 0.725 \\
\bottomrule
\end{tabular}
\end{table}

Morgan fingerprints achieved 0.900 AUROC with RF, marginally outperforming XGB (0.899). RDKit 2D descriptors achieved 0.867 AUROC, underperforming fingerprints. ChemBERTa achieved 0.802 AUROC, the lowest of the three. On balanced data, hand-engineered representations outperform learned models. This does not rule out learned representations on imbalanced data. ChemBERTa was tested on only AMES (balanced, 50% positive). Testing it on imbalanced data would clarify whether learned embeddings have fundamental limitations on class imbalance, or whether they struggle for other reasons.

\subsection{Graph Neural Network Performance Across Toxicity Tasks}

To investigate whether flexible learned representations overcome class imbalance, four GNN architectures were tested across all 17 datasets (six toxicity tasks plus twelve Tox21 endpoints). Two configurations (v1 and v2) were evaluated, differing in feature richness, model capacity, and loss function design.

Atom features (per node): atomic number, degree, formal charge, aromaticity, H-count, ring membership, hybridization state, total valence, radical electrons (9 dimensions per atom). V2's richer encoding—compared to v1's 4 features—explicitly captures properties critical for toxicity (ring membership, hydrogen bonding, valence reactivity).

Four architectures were evaluated, each with 2 convolutional layers, global aggregation (mean pooling for GCN/GAT/GraphSAGE, sum pooling for GIN), and a linear output head: GCN (degree-normalized neighbourhood mean), GIN (learnable $\varepsilon$-weighted sum, theoretically as expressive as Weisfeiler-Lehman test), GAT (multi-head attention over neighbours, learns importance weights per edge), and GraphSAGE (sample-and-aggregate with concatenation, designed for inductive generalization).

\begin{table}[H]
\centering
\caption{GNN configurations: v1 (baseline) vs v2 (improved).}
\label{tab:gnn_config}
\begin{tabular}{llrr}
\toprule
\textbf{Component} & \textbf{Description} & \textbf{v1} & \textbf{v2} \\
\midrule
Atom features & \multicolumn{1}{l}{atomic\_num, degree, charge, aromatic, \ldots} & 4 & 9 \\
Hidden dimension & GNN layer width & 16 & 64 \\
Learning rate & Adam optimizer & 0.01 & 0.001 \\
Loss function & BCE for classification & unweighted & pos\_weight \\
Datasets tested & toxicity + Tox21 endpoints & 17 & 17 \\
Models trained & (4 architectures) & 68 & 68 \\
\bottomrule
\end{tabular}
\end{table}

v1 uses minimal features and small hidden dimension, underfitting especially on regression and severe imbalance. v2 expands the atom feature space (adding H-count, ring membership, hybridization, valence, radical electrons) and capacity (hidden=64). It also applies class-weighted loss: the loss function penalizes missing a toxic molecule more heavily than a false alarm. Specifically, pos_weight = (number of non-toxic) / (number of toxic). On ClinTox with 8\% toxic, this multiplier is about 11, forcing the model to find the rare toxic examples.

GNN v2 was trained on all six toxicity datasets plus twelve Tox21 endpoints (68 total runs: 4 models $\times$ 17 datasets).

On balanced tasks (AMES, hERG), GIN achieves 0.79--0.88 AUROC, competitive with fingerprints (0.91--0.92 AUROC). Fingerprints match or exceed GNNs on all four architectures.

On severely imbalanced tasks (ClinTox 8\%, Tox21 3--15\%), fingerprints achieve 0.79 AUROC on ClinTox while GNNs achieve 0.54 AUROC with class-weighted loss. Both plateau despite different representations. At fixed 0.5 threshold, GNNs produce F1=0 because the model ranks molecules by risk (high AUROC) but never predicts "toxic." Calibrated thresholds (set to class prior: 0.08 for ClinTox) recover F1 substantially without retraining, confirming the models learn meaningful signal but need post-hoc adjustment.

On regression (LD50), fingerprints achieve $R^2 = 0.60$ while GNNs achieve $R^2 = 0.49$. Fingerprints outperform despite GNNs being theoretically better at encoding continuous properties. Graphs do encode continuous information (valence, ring geometry, hydrogen bonding) that binary fingerprints cannot: 49\% variance explained versus fingerprints' $<10\%$ on raw bit vectors shows graphs capture structure-property relationships. Yet classical ML generalizes better on this task.

Architecture ranking: GIN dominates over GCN, GAT, and GraphSAGE across all tasks. GIN wins on 13 of 16 classification datasets and on regression. But even the best GNN does not beat classical ML on imbalanced classification or regression.

% ─────────────────────────────────────────────────────────────────────────────
\section{Discussion}

\subsection{Class Imbalance as the Primary Constraint}

Three parallel experimental investigations tested different representation strategies (fingerprints, ChemBERTa, GNNs) on different dataset combinations (baseline on 6 toxicity datasets, representation comparison on AMES only, GNN analysis on all 17 datasets). Despite these differences, all converge on a single finding: class imbalance, not representation choice, determines performance on imbalanced toxicity tasks.

ClinTox provides the clearest example. With 92 toxic molecules out of 1,130 total (8\% positive), fingerprints achieve 0.79 AUROC while GNNs achieve 0.54 AUROC. Both plateau despite fundamentally different representations. The problem is not algorithmic design; it is data scarcity. Eighty-two toxic molecules cannot reliably teach any representation system to distinguish the toxic minority from the non-toxic majority with high confidence.

This challenges the common hypothesis that learned representations overcome data limitations. Instead, flexible models (GNNs, ChemBERTa) may struggle more on imbalanced data because their additional capacity allows them to fit noise in the small positive class rather than learning true toxicity signals.

\subsection{Gaps in Experimental Coverage}

The three parallel investigations created several evaluation asymmetries:

\textbf{Unequal evaluation protocol}: Fingerprints were evaluated with fixed 0.5 threshold; GNNs received post-hoc calibrated thresholds. Without consistent protocol across representations, direct comparison is incomplete.

\textbf{Different dataset coverage}: Baseline tested all 6 toxicity datasets; representation comparison tested only AMES; GNNs tested all 17 tasks. This prevents direct side-by-side comparison across the same data.

\textbf{Different imbalance strategies}: Fingerprints used rejection sampling; GNNs used class-weighted loss. Neither experiment combined both strategies on the same representation, so their synergistic effects remain unexplored.

\textbf{ChemBERTa incompletely tested}: Only evaluated on AMES (balanced, 50% positive). Whether learned embeddings fail on imbalance like GNNs remains unclear.

% ─────────────────────────────────────────────────────────────────────────────
\section{Conclusions and Future Work}

\subsection{Key Findings}

Three experimental streams (classical ML baseline, representation exploration, GNN investigation) independently arrive at the same conclusion: class imbalance is the primary performance constraint on drug toxicity prediction, regardless of representation choice. Fingerprints achieve 0.79--0.92 AUROC on balanced tasks and plateau at 0.79 AUROC on severe imbalance (ClinTox 8\%). GNNs achieve 0.79--0.88 AUROC on balanced tasks but only 0.54 AUROC on ClinTox. Learned embeddings (ChemBERTa) underperform fingerprints on balanced AMES (0.802 vs 0.900 AUROC).

\subsection{Primary Investigation: Next Two Weeks}

The team commits to one focused experiment to validate whether imbalance is fundamentally limiting or whether combined strategies can overcome it.

Apply rejection sampling (the data-side intervention that improved classical ML) to GIN on ClinTox. This directly tests whether combining data-side (rejection sampling) and model-side (class-weighted loss) interventions breaks the imbalance ceiling.

Hypothesis: If learned graph features are sound, rejection sampling should help GNNs reach at least 0.79 AUROC on ClinTox, matching fingerprints. If GNNs remain below 0.60 AUROC, the bottleneck is architectural. Timeline: one week, leaving one week for optional follow-up. Success criteria: clear yes/no answer to whether combined strategies overcome the imbalance ceiling.

\subsection{Extended Directions: Future Work}

If the primary phase succeeds, the following extensions would strengthen the work:

\textbf{3D Geometric Deep Learning}: Current GNNs use only 2D connectivity. Models like SchNet, NequIP, and Allegro incorporate 3D coordinates. Testing whether 3D geometry encodes toxicity signals that 2D misses would be significant, particularly for LD50 where acute toxicity depends on 3D geometry.

\textbf{Unified Evaluation Protocol}: Re-evaluate classical ML and GNN on all six toxicity datasets with identical conditions (scaffold splits, 5-fold CV, calibrated thresholds). This removes evaluation confounds and cleanly answers whether representation choice matters when everything else is equal.

\textbf{ChemBERTa on All Datasets}: Test pre-trained embeddings on all six toxicity datasets and Tox21. If ChemBERTa fails on imbalance like GNNs, this shows learned representations (embeddings, graphs) fundamentally struggle with data scarcity.

\textbf{Protein-Interaction Models}: Extend to DAVIS and KIBA (drug-target interaction) to test whether the imbalance ceiling generalizes beyond drug-only toxicity to broader molecular prediction tasks.

These extensions are ambitious and optional, representing the team's longer-term vision if the foundation experiment succeeds.

% ─────────────────────────────────────────────────────────────────────────────
\appendix

\section{Software and Implementation}

The project uses standard open-source libraries: RDKit for SMILES parsing and molecular graph construction,
PyTorch and PyTorch Geometric for GNN training, scikit-learn and XGBoost for baseline methods. All notebooks
and code are available in the project repository for reproducibility.

\begin{thebibliography}{20}

\bibitem{chemrev}
D.~Rogers and M.~Hahn,
``Modeling using random forests and its application to medicinal chemistry,''
\textit{Chemical Reviews}, 110(5): 3432--3465, 2010.

\bibitem{xu2019}
K.~Xu, W.~Hu, J.~Leskovec, and S.~Jegelka,
``How powerful are graph neural networks?,''
\textit{International Conference on Learning Representations (ICLR)}, 2019.

\bibitem{kipf2017}
T.~N. Kipf and M.~Welling,
``Semi-supervised classification with graph convolutional networks,''
\textit{International Conference on Learning Representations (ICLR)}, 2017.

\bibitem{velickovic2018}
P.~Velicki\'{c}, G.~Cucurull, A.~Casanova, A.~Romero, P.~Li\`{o}, and Y.~Bengio,
``Graph attention networks,''
\textit{International Conference on Learning Representations (ICLR)}, 2018.

\bibitem{hamilton2017}
W.~L. Hamilton, Z.~Ying, and J.~Leskovec,
``Inductive representation learning on large graphs,''
\textit{Advances in Neural Information Processing Systems (NeurIPS)}, 2017.

\bibitem{huang2021}
K.~Huang, T.~Fu, W.~Gao, et al.,
``Therapeutics Data Commons: Machine learning datasets and benchmarks for therapeutic discovery,''
\textit{Advances in Neural Information Processing Systems (NeurIPS), Datasets and Benchmarks Track}, 2021.

\bibitem{chawla2002}
N.~V. Chawla, K.~W. Bowyer, L.~O. Hall, and W.~P. Kegelmeyer,
``SMOTE: synthetic minority over-sampling technique,''
\textit{Journal of Artificial Intelligence Research}, 16: 321--357, 2002.

\bibitem{lin2017}
T.-Y. Lin, P.~Goyal, R.~B. Girshick, K.~He, and P.~Doll\'{a}r,
``Focal loss for dense object detection,''
\textit{IEEE International Conference on Computer Vision (ICCV)}, 2017.

\bibitem{fey2019}
M.~Fey and J.~E. Lenssen,
``Fast graph representation learning with PyTorch Geometric,''
\textit{ICLR Workshop on Representation Learning on Graphs and Manifolds}, 2019.

\bibitem{morgan1965}
A.~M. Morgan,
``The generation of a unique machine description for chemical structures: A technique developed at Chemical Abstracts Service,''
\textit{Journal of Chemical Documentation}, 5(2): 107--113, 1965.

\end{thebibliography}

\end{document}
